\textbf{Machine status: the general identities and worked example compile in the pinned Lean/Mathlib environment. Human correspondence status: awaiting Steven's review.} Build and axiom-audit evidence, source versions, and declaration coverage are recorded separately in the Chapter 0 audit. Compilation and an AI mathematical cross-check do not constitute Steven's approval.

The companion uses Mathlib's complex $L^2(X,\mu)$ and $L^\infty(X,\mu)$ spaces, Bochner integrals, continuous linear maps, orthogonal projection, and endomorphism trace. Mathlib's inner product is linear in its second argument; the shared book inner product reverses the arguments, so it is linear in the first argument as in the manuscript. A period is precisely this inner product, with a checked equality to its defining integral. Finite index types include the empty family.

For the \hyperref[finite-relative-kernel-identity]{relative kernel identity}, the formal statement proves integrability on $(X\times X,\mu\times\mu)$ and identifies the actual iterated integral with both $\langle T_Aq,p\rangle_{L^2(X,\mu)}$ and the period sum. The proof uses Mathlib's $L^2$ product-integrability theorem and the product-integral/Fubini interface. It factors each summand into its two one-variable integrals, then uses conjugate symmetry to obtain $\overline{P_q(u_j)}$. The operator pairing is checked by expanding the already constructed continuous linear operator; this is a different final step from substituting its almost-everywhere kernel representation in the prose. The integral representation itself is checked in Section 0.2.

The relative identity does not require orthonormality: the finite kernel, its operator, and the period calculation have the same meaning for any finite family in $L^2(X,\mu)$. This is a stronger statement whose orthonormal specialization is the manuscript's proposition. The relative integral theorems retain the manuscript's $\sigma$-finiteness assumption. The \hyperref[finite-relative-diagonal-identity]{relative diagonal corollary} is obtained by using an actual diagonal matrix.

For the \hyperref[finite-weighted-kernel-identity]{weighted kernel identity}, multiplication by $w$ is Mathlib's bounded H\"older multiplication map from $L^2(X,\mu)$ to itself. A checked almost-everywhere formula identifies its value with the pointwise product $wf$. The compressed endomorphism is the actual composition $P M_wT_A|_V$, where $P$ is Mathlib's orthogonal projection onto the finite span. Its trace is Mathlib's linear-map trace on that span; it is not defined to be the desired integral or diagonal sum.

The weighted proof constructs the span basis from the orthonormal family, identifies its coordinates with the book inner products, and uses the orthogonal-projection identity to compute the diagonal coefficients. A separate integrability proof justifies integrating the finite weighted kernel term by term. Orthonormality is needed for the trace identification. The weighted integral and \hyperref[finite-weighted-diagonal-identity]{weighted diagonal corollary} hold for arbitrary finite $L^2(X,\mu)$ families; the latter uses the pointwise complex modulus squared, as in the manuscript. These single-variable results do not need $\sigma$-finiteness, so their statements omit that unused assumption.

The initial formulas use the representatives supplied by Mathlib's $L^p$ classes. Separate theorems permit replacing each basis representative, each relative test function, and the weight by an almost-everywhere equal representative. Relative pairings use equality of kernels almost everywhere for the product measure. Weighted diagonal integrals use equality of the diagonals almost everywhere for the base measure, obtained from the finitely many basis-representative equalities. No weighted trace conclusion is inferred from arbitrary product-almost-everywhere equality of kernels.

The \hyperref[finite-three-measurements-example]{two-point example} uses explicit complex matrices indexed by $\operatorname{Fin}(2)$; Lean's indices $0,1$ correspond to the manuscript's points $1,2$. Counting measure produces finite sums without a normalization factor. Checked bridges identify the point-indicator kernel with the actual matrix entries and identify each ordinary, relative, and weighted measurement with its actual counting-measure integral. The example checks the two matrix-vector outputs, the weighted matrix products, both ordinary traces $3$, the relative values $0$ and $1$, and the common weighted value $w(1)+2w(2)$, including singleton weights giving $1$ and $2$. The indicator periods are checked by direct integration. The point indicators also form an orthonormal family spanning the genuine counting-measure $L^2$ space, so $V=H$ and its orthogonal projection is the identity.

The shared matrix lemmas additionally show that indicator tests recover any chosen matrix entry, that weight $1$ recovers the ordinary trace, and that matrices with the same diagonal have equal weighted measurements for every weight. The general compression proof applies to proper subspaces as well as to this example. No video script or production approval follows from these machine checks.
